\section{The uniform annulus for large negative parameters}
\label{sec:annulus}

We adapt the annulus and coefficient-separation argument of
C412~\cite[Sections 3--4]{c412}. The scaling and the threshold must
be checked anew: the centered recurrence here has a factor of two
on its neighbor sum.

For the rest of this section assume $a\le-146$ and put
\begin{equation}
 v_i=x_i+\frac12,\qquad c=-2a-\frac74,
 \qquad N=-2a-2.
 \label{eq:center}
\end{equation}
Then \eqref{eq:recurrence} becomes
\begin{equation}
 v_i^2-c=2(v_{i-1}+v_{i+1}),\qquad c=N+\frac14.
 \label{eq:center_recurrence}
\end{equation}
There is a unique integer $k\ge0$ such that
$k^2\le N\le k^2+2k$; these intervals partition the nonnegative
integers. Set
\begin{equation}
 \rho=k+\frac12,\qquad s=N-k(k+1),\qquad c=\rho^2+s.
 \label{eq:rho_s}
\end{equation}
Our hypothesis gives $N\ge290$, $k\ge17$, and
\begin{equation}
 \rho\ge\frac{35}{2},\qquad
 -\rho+\frac12\le s\le\rho-\frac12.
 \label{eq:s_range}
\end{equation}

\begin{lemma}[Five offsets about two centers]
\label{lem:annulus}
Every rational periodic orbit at $a\le-146$ satisfies
\begin{equation}
 \rho-2\le |v_i|\le\rho+2.
 \label{eq:annulus}
\end{equation}
Consequently it has a unique representation
\begin{equation}
 v_i=\varepsilon_i\rho+\delta_i,
 \qquad \varepsilon_i\in\{-1,1\},\quad
 \delta_i\in\{-2,-1,0,1,2\}.
 \label{eq:ten_symbols}
\end{equation}
\end{lemma}
\begin{proof}
Let $R=\max_i|v_i|$. The recurrence gives
\begin{equation}
 R\le2+\sqrt{c+4},\qquad |v_i^2-c|\le4R.
 \label{eq:center_bounds}
\end{equation}
The first inequality follows from $R^2-c\le4R$.
By \eqref{eq:s_range},
\[
 c+4\le\rho^2+\rho+\frac72<(\rho+1)^2,
\]
where the strict inequality holds for $\rho>5/2$.
Thus $R<\rho+3$. By Lemma~\ref{lem:integrality}, both $R$ and
$\rho$ are half-integers, so $R\le\rho+2$.

If $|v_i|\le\rho-3$ at an index, then
\[
 c-v_i^2\ge\rho^2-\rho+\frac12-(\rho-3)^2
 =5\rho-\frac{17}{2}>4\rho+8\ge4R.
\]
The strict inequality requires $\rho>33/2$, which our threshold
satisfies. This contradicts \eqref{eq:center_bounds}.
The half-integral lattice therefore gives the lower bound in
\eqref{eq:annulus}. The two center intervals are disjoint and
their offsets are integers, proving uniqueness in
\eqref{eq:ten_symbols}.
\end{proof}

\begin{lemma}[Exact separation into six symbols]
\label{lem:separation}
In \eqref{eq:ten_symbols}, every offset is even and $s$ is divisible
by four. Writing $\delta_i=2d_i$ and $t=s/4$, the entire cyclic
sequence satisfies exactly
\begin{equation}
 \varepsilon_{i-1}+\varepsilon_{i+1}=2\varepsilon_i d_i,
 \qquad d_{i-1}+d_{i+1}=d_i^2-t,
 \qquad d_i\in\{-1,0,1\}.
 \label{eq:local}
\end{equation}
\end{lemma}
\begin{proof}
Substitution into \eqref{eq:center_recurrence} gives
\begin{equation}
 2\rho(\varepsilon_{i-1}+\varepsilon_{i+1}
              -\varepsilon_i\delta_i)
 =\delta_i^2-2\delta_{i-1}-2\delta_{i+1}-s.
 \label{eq:separation}
\end{equation}
The right side has absolute value at most
\[
 12+|s|\le\rho+\frac{23}{2}<2\rho.
\]
The coefficient multiplying $2\rho$ on the left is an integer.
It must be zero, and so must the right side. The sum of two
neighboring signs is even, hence every $\delta_i$ is even.
On writing $\delta_i=2d_i$, the vanishing right side gives
$s=4(d_i^2-d_{i-1}-d_{i+1})$. Division gives
\eqref{eq:local}.
\end{proof}

The six pairs $(\varepsilon_i,d_i)$ and both local equations
\eqref{eq:local} are exactly the symbol system already classified
in C412. We reproduce its elementary word exhaustion to make the
dependence transparent and to keep the present proof readable
without an external proof-file lookup.
